% Answers to Chapter 14, from lectures/L14/appendix/c_solutions.md; the self-test from
% lectures/L14/README.md.

\facitavsnitt{c14}{Kapitel 14}{Integraler}
\begin{facit}

\svar{1.1} \sv{a} $f_1'(x) = 2xe^{4x}(2x + 1)$
\sv{b} $f_2'(x) = x^2(3\ln 2x + 1)$
\sv{c} $f_3'(x) = \frac{3x^2\ln x^2 - 2x^2}{(\ln x^2)^2}
= \frac{x^2(3\ln\abs{x} - 1)}{2(\ln\abs{x})^2}$

\svar{2.1} \sv{a} $p(t) = 5t\,\text{W}$ \sv{b} $P(t) = \frac{5t^2}{2} + C$
\sv{c} $w(t) = \frac{5t^2}{2} + C$, samma uttryck som $P(t)$.
\sv{d} $C = 0$, så $w(t) = \frac{5t^2}{2}\,\text{J}$.
\sv{e} $22{,}5\,\text{J}$

\svar{3.1} \sv{a} $\frac{x^4}{4} + C$ \sv{b} $2x^2 + C$ \sv{c} $5x + C$ \sv{d} $x^2 + 3x + C$
\sv{e} $\frac{x^3}{3} - 2x^2 + x + C$

\svar{3.2} \sv{a} $e^x + C$ \sv{b} $\frac{e^{3x}}{3} + C$ \sv{c} $\ln\abs{x} + C$
\sv{d} $-\cos(x) + C$ \sv{e} $\frac{\sin(2x)}{2} + C$ \sv{f} $-\frac{\cos(4x)}{4} + C$

\svar{3.3} \sv{a} $F'(x) = \frac{3x^2}{3} = x^2 = f(x)$
\sv{b} $F'(x) = -\frac{1}{2} \cdot \bigl(-2\sin(2x)\bigr) = \sin(2x) = f(x)$
\sv{c} $F'(t) = -2 \cdot \bigl(-\frac{1}{2}\bigr)e^{-t/2} = e^{-t/2} = f(t)$
\sv{d} Derivatan av en konstant är noll, så $F(x) + C$ har samma derivata för varje värde på $C$.

\svar{3.4} \sv{a} $8$ \sv{b} $8$ \sv{c} $12$ \sv{d} $\frac{2}{3}$ \sv{e} $9$

\svar{3.5} \sv{a} $e - 1 \approx 1{,}72$ \sv{b} $2$ \sv{c} $1$ \sv{d} $1$
\sv{e} $1 - e^{-2} \approx 0{,}865$

\svar{3.6} \sv{a} $8$ \sv{b} En triangel med basen $4$ och höjden $4$: $\frac{4 \cdot 4}{2} = 8$.
\sv{c} $12$, som rektangeln med bredden $4$ och höjden $3$. \sv{d} $\frac{8}{3} \approx 2{,}67$

\svar{3.7} \sv{a} $-2$ \sv{b} $2$ \sv{c} $0$
\sv{d} Integralen räknar area med tecken: delen under $x$-axeln och den lika stora delen ovanför
tar ut varandra. Den totala arean är $\abs{-2} + \abs{2} = 4$.

\svar{3.8} \sv{a} $f(x) = 2x^2 - 3x + C$ \sv{b} $C = 5$ \sv{c} $f(x) = 2x^2 - 3x + 5$ \sv{d} $7$

\svar{3.9} \sv{a} $u(t) = 3t^2 + C$ \sv{b} $C = 2\,\text{V}$ \sv{c} $u(t) = (3t^2 + 2)\,\text{V}$
\sv{d} $50\,\text{V}$

\svar{3.10} \sv{a} $q(t) = (2t^2 + 2t)\,\text{C}$ \sv{b} $24\,\text{C}$ \sv{c} $60\,\text{C}$
\sv{d} $36\,\text{C}$ \sv{e} $\Bigl[2t^2 + 2t\Bigr]_3^5 = 60 - 24 = 36\,\text{C}$, samma svar.

\svar{3.11} \sv{a} $-\frac{2\cos(100\pi t)}{100\pi} + C$
\sv{b} $\frac{4}{100\pi} \approx 0{,}0127\,\text{C}$
\sv{c} $0$
\sv{d} Strömmen är positiv under den första halvperioden och lika stor men negativ under den
andra, så netto passerar ingen laddning under en hel period.

\svar{3.12} \sv{a} $i'(t) = 3\,\text{A/s}$ \sv{b} $p(t) = 3{,}6t\,\text{W}$
\sv{c} $w(t) = 1{,}8t^2\,\text{J}$ \sv{d} $7{,}2\,\text{J}$
\sv{e} $i(2) = 6\,\text{A}$ ger $w = \frac{0{,}4 \cdot 6^2}{2} = 7{,}2\,\text{J}$, samma svar.

\svar{3.13} \sv{a} $w(t) = (5t + t^2)\,\text{J}$ \sv{b} $150\,\text{J}$ \sv{c} $100\,\text{J}$
\sv{d} $15\,\text{W}$

\svar{3.14} \sv{a} $3$ \sv{b} $4$ \sv{c} $\frac{20}{\pi} \approx 6{,}37\,\text{A}$
\sv{d} $\frac{2}{\pi} \approx 0{,}637$, alltså cirka $64\,\%$ av amplituden.

\svar{3.15} \sv{a} $8$ \sv{b} $10$ \sv{c} $9$ \sv{d} $-5$

\svar{3.16} \sv{a} $F'(x) = 3x^2 - 2$ \sv{b} $x^3 - 2x + C$
\sv{c} En konstantterm försvinner vid derivering, eftersom dess derivata är noll, så vid
integreringen går det inte att veta vilken konstant som fanns och man skriver ett allmänt $C$.
\sv{d} $5x^4$

\svar{3.17} \sv{a} $a = 5$ \sv{b} $a = 4$ \sv{c} $k = 5$ \sv{d} $4\,\text{s}$

\svar{3.18} \sv{a} Man har deriverat i stället för att integrera: $\frac{x^3}{3} + C$.
\sv{b} Den inre faktorn ska hamna i nämnaren: $\frac{e^{2x}}{2} + C$.
\sv{c} Tecknet är fel: $-\cos(x) + C$.
\sv{d} Den undre gränsen är bortglömd: $8 - 1 = 7$.
\sv{e} Falskt. Konstanten tar ut sig själv: $\bigl(F(b) + C\bigr) - \bigl(F(a) + C\bigr) =
F(b) - F(a)$.
\sv{f} Absolutbeloppet saknas; $\ln x$ är odefinierat för $x < 0$. Rätt är $\ln\abs{x} + C$.

\svar[Testa dig själv]{T} \sv{1} $10$ \sv{2} $F(x) = 2x^2 + e^x + C$

\end{facit}
